\documentclass[11pt]{article}

% "review" anonymizes and numbers lines; "preprint" for a non-anonymous version.
\usepackage[review]{acl}

\usepackage{times}
\usepackage{latexsym}
\usepackage[T1]{fontenc}
\usepackage[utf8]{inputenc}
\usepackage{microtype}
\IfFileExists{inconsolata.sty}{\usepackage{inconsolata}}{}
\usepackage{newunicodechar}
\newunicodechar{∀}{\ensuremath{\forall}}
\newunicodechar{→}{\ensuremath{\to}}
\usepackage{booktabs}
\usepackage{array}
\newcolumntype{L}[1]{>{\raggedright\arraybackslash}p{#1}}
\usepackage{amsmath,amssymb}

\newcommand{\yes}{\textit{yes}}
\newcommand{\no}{\textit{no}}
\newcommand{\unk}{\textit{unknown}}
\newcommand{\nm}[1]{\texttt{#1}}

\title{Every ``Yes'' Is a Proof: Where a Logic-Based NLI Pipeline Goes Wrong,\\and How It Can Refuse Instead}

\author{Anonymous}

\begin{document}
\maketitle

\begin{abstract}
When every verdict of a natural language inference system is a
machine-checked proof, the prover cannot be the source of a wrong
answer: every error is an error of translation. We build such a
system, from FraCaS's own English to Lean~4 propositions closed by a
fixed, unattended prover, and evaluate it by its wrong answers. Two
design rules carry the weight. An ambiguous sentence is translated in
each of its readings, and the system answers only if every reading
agrees. A fused predicate name is refused when it would stand for two
different phrases: we show that this one collision is the mechanism
behind every naming error the pipeline has made, and that a check for
it subsumes two earlier safeguards. On the 259 FraCaS problems the
pipeline reaches (75.7\%), it matches 240 with no wrong answer;
this is a development-set figure. On held-out data it is a precision
system with very low recall. On MED, with the pipeline frozen, it gave
3 wrong answers in 140; one was a real translation error, which the
name check now refuses, and two are defects of how MED's pairs were
built. On SICK's test split, pre-registered after the fix, it
answered 52 of 4,906 pairs, all correctly (95\% lower bound on
precision 93.2\%), for a recall of 2.5\%.
\end{abstract}

\section{Introduction}

Logic-based approaches to natural language inference (NLI) translate
premises and hypothesis into a logic and ask a prover whether the
hypothesis follows \citep{mineshima-etal-2015-higher,abzianidze-2015-tableau,bernardy-chatzikyriakidis-2017-type}.
Their appeal is that an answer comes with a reason. But a proof is
only as good as the formula it proves. If the prover is sound and its
proofs are checked, then a wrong \yes{} or \no{} can come from one
place only: the translation said something the English does not.

This paper takes that observation as a design principle and an
evaluation criterion. We build a pipeline from English to Lean~4
\citep{demoura-ullrich-2021-lean4} whose every \yes{} and \no{} is a
kernel-checked proof, and we judge it first by its wrong answers.
Accuracy matters too, but a system of this kind can trade coverage for
soundness in a way a classifier cannot: it can decline to translate
what it might translate wrongly, and say \unk{}.

The pipeline parses FraCaS's original English
\citep{cooper-etal-1996-fracas} with a restricted Grammatical Framework
grammar \citep{ranta-2011-gf}, ranks the readings generically, resolves
ellipsis and anaphora, translates, adds problem-independent background
knowledge, and runs a fixed tactic chain (\S\ref{sec:pipeline}). Three
contributions follow from asking where its wrong answers come from.

\begin{enumerate}
\item \textbf{The every-reading rule} (\S\ref{sec:readings}). Six kinds
of ambiguity are translated in a weak and a strong reading, and the
system answers only if every consistent assignment of readings proves
the same thing. Ablations show that committing to either reading, or
answering if \emph{any} reading proves the hypothesis, gives between
one and six wrong answers on FraCaS, and buys at most one problem.
\item \textbf{The name check} (\S\ref{sec:names}). Like other
pipelines, ours names a predicate after the content words of its
phrase (``spend a lot of time'' becomes a predicate
\nm{spendALotOfTime}). Names are bound variables of the claim, so a
name that leaves a word out is harmless by itself. It is wrong when the
same name also names a \emph{different} phrase. We show that this
collision is the mechanism behind every naming error the pipeline has
made, on FraCaS and on held-out data, and refuse it with one check.
\item \textbf{Held-out tests with a frozen pipeline} (\S\ref{sec:heldout}).
On MED \citep{yanaka-etal-2019-neural}, run before the name check, the
pipeline gave three wrong answers in 140; we analyse each. On SICK
\citep{marelli-etal-2014-sick}, run after it with a protocol committed
before any claim was generated, it gave none in 52, at a recall of 2.5\%.
\end{enumerate}
The code, the generated claims and the verdicts of every run cited here
will be released with the paper.

We state the main limitation at the outset. Every FraCaS figure here is
a development-set figure: each construction, reading, background rule
and prover step was added after seeing a FraCaS problem fail. The
grammar is a controlled language fit to FraCaS, and it reaches 18.5\% of
MED and 28.9\% of SICK. This is a paper about precision, not coverage.

\section{The pipeline}
\label{sec:pipeline}

\paragraph{Parsing and ranking.} Each sentence is tokenized by one
generic rule and parsed with a restricted GF grammar. Parsing with the
full grammar is not usable: one Section~1 sentence gave 69,115 readings
in five minutes. The restricted grammar keeps only the functions and
words used by hand-built trees of FraCaS sentences, plus constructions
FraCaS's own wording needs (relative ``who'', existential ``there
are'', VP ellipsis, gapping, fragments, the progressive), each confirmed
by a linearization check. It does not say which tree goes with which
sentence. Readings are ranked by size, with ties broken by attaching
adverbials as low as possible. On the project's 486 sentences with a confirmed hand-built
tree (its own versions of FraCaS sentences, some reworded), every hand
tree is among the readings; for the 406 the translator covers, the
top-ranked reading translates to the same proposition as the hand tree
for 405. The tie-breaker was chosen on these same sentences.

\paragraph{Discourse.} Ellipsis and one-anaphora are resolved by copying
the antecedent's verb phrase or noun; a gapped clause copies its
antecedent with each remnant replacing its parallel adverbial. Where the
antecedent is ambiguous (``John said Mary wrote a report, and Bill did
too''), the problem is not translated. Possessives bound to a quantified
subject are functions of their possessor, and pronouns after ``NAME
V-ed a N'' are E-type descriptions of their antecedent.

\paragraph{Translation and assembly.} Each tree is translated to a Lean
proposition over generalized quantifiers (\nm{Every}, \nm{Some},
\nm{AtLeast}, \nm{Exactly}, \nm{More}, \ldots). Nouns, verb phrases and
other predicates become predicate variables with fused names. One
generic rule assembles a claim: bind every free name universally over
an arbitrary finite domain, and chain the premises into the hypothesis
(a \yes{} claim) or into its negation (a \no{} claim). There are no
per-problem cases. Figure~\ref{fig:claim} shows one claim.

\begin{figure}[t]
\small
\begin{verbatim}
∀ {A : Type} [Fintype A]
  (britishDelegate obtainResult delegate
     : A → Prop)
  (Many : (A → Prop) → (A → Prop) → Prop),
  (∀ P Q, Many P Q → Some P Q) →
  (∀ P Q R, Many P Q → (∀ x, Q x → R x)
            → Many P R) →
  (∀ x, britishDelegate x → delegate x) →
  Many britishDelegate obtainResult →
  Many delegate obtainResult
\end{verbatim}
\caption{FraCaS 56 (``Many British delegates obtained interesting
results from the survey'' $\Rightarrow$ ``Many delegates \ldots'', gold
\unk{}), its \yes{} claim on the proportional reading of ``many'', with
one background premise (a name shortened, instance arguments omitted). ``Many'' is opaque:
the claim states only properties both readings share. The prover does
not close this claim, so the answer is \unk{}; on the cardinal reading
(``many'' as \nm{Some}) it would close.}
\label{fig:claim}
\end{figure}

\paragraph{Background knowledge.} Each problem is scored twice: as
translated, and with problem-independent background knowledge added as
premises. Most of it is derived from the parse trees: deleting an
intersective modifier in an upward-monotone position (never a
non-subsective adjective like ``former'', never under \nm{every} or
\nm{no} objects, never a durative ``for two hours'' from a verb phrase
that culminates). The rest is a small closed lexicon: antonym contraries
over the same noun, the adjective classes of
\citet{kamp-partee-1995-prototype}, relations between time adverbials,
and the event ``finish NP'' stands for \citep{pustejovsky-1995-generative}.
Nothing in it is keyed by problem or gold answer, but it was written
with FraCaS in view.

\paragraph{Proving.} A fixed chain of tactics is tried on each claim,
each under a fresh heartbeat budget: a certified monotonicity
fragment, then \nm{aesop}, \nm{grind}, \nm{omega} and counting
arguments over cardinal quantifiers. A \yes{} answer means every
reading's \yes{} claim closed; a \no{} answer means every reading's
\no{} claim closed; a problem where both close counts as a wrong answer;
anything else is \unk{}. Every closed claim is checked by Lean's kernel.
Continuous integration regenerates all claims from the English and fails
on any difference from the committed copy, so the scored claims are
reproducible.

\paragraph{Refusing.} Where a step cannot translate a sentence
faithfully, the problem is not translated, and counts as \unk{}. Besides
the name check (\S\ref{sec:names}), the pipeline refuses ellipsis with
two possible antecedents, a coordinated subject whose collective reading
is live (``Smith and Jones signed two contracts''), an object pronoun
with no resolved antecedent, and sentential negation except with a
proper-name subject, where its scope is unambiguous.

\paragraph{Where errors can come from.} Since the kernel checks every
proof, a wrong answer means a claim is provable that the English does
not support. That can be the parse, a reading, the translator, a name,
the assembly, or a background premise. All of these are
\emph{translation} in the sense of this paper: everything between the
English and the claim.

\section{Ambiguity: the every-reading rule}
\label{sec:readings}

An earlier version of the pipeline gave five wrong answers on FraCaS,
all from translation shortcuts that committed to one reading: ``the
inhabitants voted'' read distributively, ``many'' as ``some'', ``both''
as ``every'', plural ``some accountants'' as ``at least one''. The fix
was not to choose better. Six kinds of ambiguity are now translated in
a weak and a strong reading: a definite plural subject (collective or
distributive), a bare plural with a relative clause (existential or
generic), ``many'' (cardinal or proportional), plural ``some Ns''
(number-neutral or at least two), a bare numeral (at least or exactly),
and ``few'' (a small number or a small proportion). A proportional
reading is an opaque relation stated only with the properties every
reading shares (Figure~\ref{fig:claim}). Each problem is translated
under every \emph{consistent} assignment of readings, the same word read
the same way in all its sentences, and answered only if every assignment
agrees.

FraCaS sometimes commits to one reading itself. Problems 107 and 109
are the same English (``Just one accountant attended the meeting''
$\Rightarrow$ ``Some accountants attended the meeting'') with gold
answers \yes{} and \no{}: the number-neutral and at-least-two readings of
``some accountants''. No consistent system answers both; the
every-reading rule answers \unk{} to both. The ablations
(\S\ref{sec:ablations}) measure the rule against its alternatives.

\section{Naming: the name check}
\label{sec:names}

\paragraph{Fused names.} The translator names a predicate after the
words of its phrase that belong to categories it lists, and silently
drops the rest. ``Tom said he didn't drink black coffee'' should be a
relation between Tom and a negated clause; if ``said'' and ``didn't''
are not listed, the verb phrase becomes \nm{drinkBlackCoffee}. Is a
name that drops words wrong? Not by itself. Every name is a bound
variable of its claim, interpreted arbitrarily: a claim about
\nm{drinkBlackCoffee} holds for every predicate, including ``said he
didn't drink black coffee''.

\paragraph{The collision.} A name goes wrong when it also names a
\emph{different} phrase in the same problem, because whatever is
stated of one is then stated of the other. In that example (MED~2718), the outer
verb phrase and the one inside it (``drink black coffee'') get the same
name, so a weakening that holds of the inner one (``drink black coffee''
$\Rightarrow$ ``drink coffee'') is applied to the outer one, and the
pipeline proved ``Tom said he didn't drink coffee'', gold neutral.
Before determiners were kept in names, ``spend a lot of time'' and
``spend time'' were both \nm{spendTime}, which gave FraCaS 30, 32 and
47 wrong answers (``Neither commissioner spends a lot of time at home''
$\Rightarrow$ ``Neither commissioner spends time at home'', gold \unk{}).

More precisely, let a \emph{fine} translation give each distinct phrase
its own predicate variable. The pipeline's claim is obtained from it by
identifying the variables of phrases that share a name. If no two
distinct phrases share a name (phrases that differ only in what the
deliberate omissions below allow count as one), this is a renaming of
bound variables and the two claims are equivalent. If two do, the pipeline's claim is
an instance of the fine one with two variables made equal, and can be
provable when the fine one is not. So a dropped word can cost soundness
only when it makes the naming map non-injective on the phrases of a
problem.

\paragraph{The check.} \nm{names.py} computes, for every phrase of a
category the translator names (common nouns, verb phrases, complements,
relative clauses, adjectival phrases, adverbials, read off the grammar),
its name, as the translator computes it, and the words its name leaves
out. A problem is refused when a name its claim uses names two phrases
that differ in what they leave out. What a name may leave out without
counting is a short, explicit list, each entry with its reason: articles
and grammatical number (a determiner becomes the quantifier, outside the
name), tense (the claims are tenseless), positive polarity, relative
pronouns, and a floated ``all''. A negation, a verb taking a clause, an
adverb, a pronoun or a conjunction counts. Two further cases need care.
First, background knowledge states facts about phrases that occur in no
sentence (a weakened verb phrase, the event ``finish'' stands for); the
check includes them. Second, a construction with no word of its own
counts as left out when its mark is missing from a name: the
progressive, which names mark (\nm{progBuildMtalk}), and the passive,
which they do not. The check calls the translator's own naming code, so
it also guards later changes to it. It has 13 unit tests, run in CI.

On FraCaS the check refuses one reached problem, 119 (``No student used
her workstation'' and ``Mary used a workstation'' would share
\nm{useWorkstation}), which was \unk{} before; no verdict changes. Its
value shows when another safeguard is removed (\S\ref{sec:ablations})
and on held-out data (\S\ref{sec:heldout}).

\section{Development results: FraCaS}
\label{sec:fracas}

\paragraph{Reach.} Of FraCaS's 342 problems with a hypothesis (4 of the
346 in MacCartney's version have none), 259 (75.7\%) are carried
through. Of the rest, 65 stop at a sentence the grammar does not parse,
17 at a reading the translator does not cover (including 119, refused by
the name check), and one at a name that is a Lean keyword. The remaining
parse failures are mostly hard phenomena: tense and aspect,
embedded questions, antecedent-contained ellipsis, cross-sentence
anaphora, and numbers written as digits.

\begin{table}[t]
\centering\small
\begin{tabular}{lrr}
\toprule
& matched & wrong \\
\midrule
always \unk{} (baseline) & 90 / 259 & -- \\
as translated & 181 / 259 & 0 \\
with background knowledge & 240 / 259 & 0 \\
\bottomrule
\end{tabular}
\caption{FraCaS, the 259 reached problems (development set). Counting
the 83 unreached problems as \unk{}, corpus-wide accuracy is 57.6\% and
74.9\%, against a 31.0\% baseline.}
\label{tab:fracas}
\end{table}

\paragraph{Results.} Table~\ref{tab:fracas} gives the results. No
reached problem is answered wrongly in either variant. Every mismatch is
\unk{} against a \yes{} or \no{} label. With background knowledge there
are 19: 13 where FraCaS commits to one reading of an ambiguous sentence
(13, 76, 86, 87, 100, 107, 109, 129, 160, 230, 233, 239, 250), a slip in
FraCaS itself (69), three inferences the pipeline deliberately does not
make (140: what is said closed under entailment; 199, 218: ``successful''
and ``clever'' as intersective), and two where a fronted ``since'' would
have to force a continuous reading (255, 257).

\paragraph{Against systems that parse FraCaS's English.}
\citet{mineshima-etal-2015-higher}'s ccg2lambda parses with a
statistical CCG parser and proves with Coq, on FraCaS sections 1, 2, 5,
6, 8 and 9 minus seven problems without a well-defined answer. We
score this pipeline on the same 181 problems and labels, counting the
ones it does not reach as \unk{}.
Table~\ref{tab:ccg} compares per section, with
\citet{haruta-etal-2020-logical}'s successor system where it reports a
section. The comparison favours us in one important way: our grammar is
a controlled language fit to FraCaS, while ccg2lambda's parser is
trained on newswire and its lexicon is not specific to FraCaS. Both are
development-set figures.

\begin{table}[t]
\centering\small
\setlength{\tabcolsep}{3.5pt}
\begin{tabular}{lrrrr}
\toprule
section (problems) & ours & ours+BK & ccg2l. & H20 \\
\midrule
quantifiers (74) & .76 & .96 & .77 & .92 \\
plurals (33) & .70 & .70 & .67 & -- \\
adjectives (22) & .32 & .91 & .68 & .86 \\
comparatives (31) & .48 & .65 & .48 & .77 \\
verbs (8) & .50 & .88 & .62 & -- \\
attitudes (13) & .77 & .77 & .77 & -- \\
\midrule
total (181) & .635 & .834 & .69 & -- \\
\bottomrule
\end{tabular}
\caption{Accuracy on the 181 FraCaS problems ccg2lambda is scored on
\citep[Table 3]{mineshima-etal-2015-higher}. H20:
\citet{haruta-etal-2020-logical}, on 73/22/31 problems without
hand-coded tagging rules. BK: background knowledge. Our columns count
unreached problems as \unk{}.}
\label{tab:ccg}
\end{table}

FraCoq \citep{bernardy-chatzikyriakidis-2017-type,bernardy-chatzikyriakidis-2019-wide,bernardy-chatzikyriakidis-2021-applied}
is the strongest published logic-based FraCaS system, and it is ahead
of this pipeline on all three of its problem sets. Scored under FraCoq's own amended labels
and rules, the pipeline with background knowledge gets 81.9\% of the
2017 problem set (FraCoq 0.83), 74.1\% of the 2019 set (0.89) and 73.9\%
of the 2021 set (0.81). The settings differ: FraCoq starts from a
treebank with one manually chosen tree per sentence
\citep{ljunglof-siverbo-2011-bilingual}, its proofs needed human
intervention, and it counts an ambiguous problem as \yes{} if one reading
entails the hypothesis.

\subsection{Which safeguards prevent wrong answers}
\label{sec:ablations}

Each safeguard is switched off in turn and the problems re-scored
(Table~\ref{tab:ablations}). Rows that change the reading rule reuse the
pipeline's claims, scored one reading at a time; the others regenerate
the claims with the safeguard patched out.

\begin{table}[t]
\centering\small
\setlength{\tabcolsep}{3pt}
\begin{tabular}{p{3.6cm}cc}
\toprule
switched off & wrong & wrong+BK \\
\midrule
nothing (the pipeline) & 0 & 0 \\
every reading $\to$ weak only & 1 & 2 \\
every reading $\to$ strong only & 4 & 4 \\
every reading $\to$ any reading & 5 & 6 \\
presupposition in \no{} claims & 0 & 1 \\
telicity & 0 & 2 \\
non-subsective / non-veridical exclusions & 0 & 2 \\
determiners, numerals, progressive in names & 0 & 0 \\
\quad the same, before the name check & 3 & 5 \\
the name check & 0 & 0 \\
\bottomrule
\end{tabular}
\caption{Wrong answers on FraCaS with one safeguard switched off.
``Any reading'' answers \yes{} if some reading proves the hypothesis,
the rule FraCoq's amended labels follow. Matched problems change by at most
one in the reading rows, and by 2 and 6 in the names row.}
\label{tab:ablations}
\end{table}

\paragraph{Committing to a reading buys almost nothing.} Weak readings
alone match one more problem and answer one or two wrongly (56, 109);
strong readings alone match none more and answer four wrongly; the
any-reading rule answers five or six wrongly, counting two problems
where both directions close. The every-reading rule's price is one
problem.

\paragraph{Each semantic safeguard prevents specific errors.} Telicity,
which stops ``was winning the contract'' from entailing that it was won
\citep{dowty-1979-word} and a durative ``for two hours'' from being
deleted from a verb phrase that culminates, prevents two (329, 304).
Excluding non-subsective adjectives and non-veridical adverbs from
modifier deletion prevents two (198, 201). Requiring a \no{} claim to
establish a partitive's presupposition \citep{strawson-1950-referring}
prevents one (78).

\paragraph{The name check subsumes two safeguards.} Determiners,
numerals and the progressive's mark were added to names over the
pipeline's development. Dropping all three gave 3 and 5 wrong
answers before the name check (30, 32, 47, 327, 329). With the check, the
13 problems whose names then collide are refused, and none is answered
wrongly. Likewise, without the guard against unresolved object pronouns the
check refuses the two problems that guard catches (134, 135). Switched
off on its own, the check changes nothing: 119 is reached again and
answered \unk{}.

These ablations show what each safeguard does on FraCaS. Since each was
added after a FraCaS error, they do not show that it would be needed, or
enough, on unseen problems. That is what the held-out tests are for.

\section{Held-out tests}
\label{sec:heldout}

Both tests freeze everything that decides an answer: constructions,
ranker, discourse resolution, readings, translator, name check (SICK
only), background knowledge and prover. One thing is added, by a fixed
rule: every open-class entry of the GF English dictionary for a word the
FraCaS grammar does not know, and a proper noun for a capitalized word
with none. No closed-class word, construction or lexical knowledge
(``dog'' $\sqsubseteq$ ``animal'') is added. The pipeline is not changed
in response to a test's results; figures computed after a change are
labelled post-hoc. The verdicts of both CI runs are kept with the
repository, and a script recomputes every figure from them.

\subsection{MED, before the name check}

MED \citep{yanaka-etal-2019-neural} has 5,382 single-premise problems
labelled entailment or non-entailment (neutral); we exclude the 37 taken from FraCaS. Of
the remaining 5,345, 989 are reached (18.5\%): 1,891 need a closed-class
word the grammar lacks, 1,807 do not parse, and 658 do not translate. A
\yes{} on neutral or a \no{} on entailment is wrong; \unk{} on
entailment is a miss.

\begin{table}[t]
\centering\small
\setlength{\tabcolsep}{4pt}
\begin{tabular}{lrrrr}
\toprule
& reached & proved & decided & wrong \\
\midrule
as translated & 989 & 1 & 2 & 1 \\
with BK (frozen) & 989 & 137 & 140 & 3 \\
with BK + name check$^\dagger$ & 961 & 137 & 139 & 2 \\
\bottomrule
\end{tabular}
\caption{MED (5,345 problems; 477 reached problems are entailments).
Proved: entailments answered \yes{}. $^\dagger$Post-hoc: the name check
was written in response to this run; it only removes problems, so the
remaining verdicts are the frozen run's.}
\label{tab:med}
\end{table}

Table~\ref{tab:med} gives the results. The pipeline never answered
\no{}. Background knowledge proves 136 of the 137 entailments, mostly by
deleting a modifier in an upward position, which is what MED's upward
problems test; without lexical knowledge its lexical and most of its
downward problems stay \unk{}. With background knowledge, 137 of 140
answers are correct (97.9\%). The three wrong answers, each a \yes{}
on a neutral problem:

\begin{itemize}
\item \textbf{2718} is a real translation error, the collision of
\S\ref{sec:names}: ``Tom said he didn't drink black coffee''
$\Rightarrow$ ``Tom said he didn't drink coffee''. On FraCaS, every
attitude report was handled by a dedicated reading, which declines a
negated complement or a pronoun subject; 2718 has both, so it fell
through to the default naming, which drops the verb and the negation.
\item \textbf{536} (``A cold bubbly bath refreshed him'' $\Rightarrow$
``A cold bath refreshed him'') and \textbf{2504} (premise and hypothesis
the same sentence) follow from how MED built its crowdsourced pairs.
Labels were computed from monotonicity and kept where all three
workers agreed; then, where possible, each pair was swapped and its
label reversed \citep[\S3.1.3]{yanaka-etal-2019-neural}.
536 is 535 swapped with its label left unreversed (MED has 8 such
pairs); 2504 should be 2503 swapped, but its hypothesis repeats its premise
(MED has 6 pairs with the same sentence twice).
\end{itemize}

We count all three as wrong, as the protocol says. Against labels
corrected for 536 and 2504, one answer of 140 is wrong: 2718. Either
way, ``no wrong answers'' did not survive as a property of the pipeline.

\paragraph{How far the gap reached.} Applying the name check, after the
run and without changing it, to the frozen pipeline's translations: 28
reached problems have a name standing for two phrases. In 18 the name
drops the negation of a relative clause (``every alien who did not
throw balls''), in 3 ``said'' and a negated complement, in 4
``believes'', in 2 ``even'', and in 1 ``no longer''. Only 2718 became
a wrong answer. The other 27 came out \unk{}, but some would not stay
so once lexical knowledge is added: MED~4897's claim follows from
``quarter $\sqsubseteq$ coin'', while its English (``every alien who did
not use coins is hairy'' $\Rightarrow$ ``\ldots\ did not use quarters
\ldots'') does not.

\paragraph{Against other systems.} Whole-MED accuracy is not comparable
at 18.5\% reach; NeuralLog \citep{chen-etal-2021-neurallog} reports 93.4.
\citet{haruta-etal-2020-logical} score a 498-problem subset that needs no
world knowledge, at .84. Scoring our frozen run on the subset MED's tags
give (unreached problems as \unk{}): .600, against a .584
all-\unk{} baseline, with 8 of 207 entailments proved and no wrong
answer (MED's labels have changed since: theirs split 215/283, ours
207/291). On that subset the pipeline is at the baseline.

\subsection{SICK, after the name check}

Since the name check was written in response to MED, MED cannot test the
pipeline as it now is. We ran a second test on SICK's SemEval test split
\citep{marelli-etal-2014-sick}: 4,906 pairs (1,404 entailment, 712
contradiction, 2,790 neutral), three-way labelled. The protocol,
including the scoring rules and every figure to be reported, was
committed with the data, generator and scorer before any claim was
generated, with the pipeline frozen at that commit.

\begin{table}[t]
\centering\small
\setlength{\tabcolsep}{3.5pt}
\begin{tabular}{lrrrr}
\toprule
& decided & wrong & precision & recall \\
\midrule
as translated & 19 & 0 & 100\% & 0.9\% \\
with BK & 52 & 0 & 100\% & 2.5\% \\
\midrule
LangPro (2017) & & & 98.0 & 58.1 \\
Mart\'inez-G\'omez et al. & & & 97.0 & 63.6 \\
\bottomrule
\end{tabular}
\caption{SICK test split, pre-registered (1,416 of 4,906 reached).
Recall is over all 2,116 entailment and contradiction pairs, unreached
pairs counting as \unk{}. With 52 of 52 correct, the 95\%
Clopper--Pearson lower bound on precision is 93.2\%. The last two rows
are as tabulated by \citet[Table 3]{chen-etal-2021-neurallog}, citing
\citet{abzianidze-2017-langpro} and
\citet{martinez-gomez-etal-2017-demand}, on SICK's 4,927 test pairs; the
release we use has 4,906.}
\label{tab:sick}
\end{table}

The pipeline reaches 1,416 pairs (28.9\%); 2,714 do not parse, 454
parse but do not translate, and 322 need a closed-class word.
Table~\ref{tab:sick} gives the results. It proves 49 entailments, nearly
all by deleting a modifier in an upward position (``A wild deer is
jumping a fence'' $\Rightarrow$ ``A deer is jumping a fence''), and 3
contradictions of the form ``A N is V-ing'' / ``No N is V-ing''. Of the
reached pairs, 543 of 592 entailments and 59 of 62 contradictions stay
\unk{}; SICK's entailments typically need lexical knowledge (``man''
$\sqsubseteq$ ``person'') or paraphrase. Accuracy over the whole split,
for context only, is .579 against a .569 all-\unk{} baseline.

The precision sample is small: 52 of 52 correct
cannot be told apart from LangPro's 98.0 or
\citeauthor{martinez-gomez-etal-2017-demand}'s 97.0
\citep{clopper-pearson-1934}. What the test shows is narrower: on data
nothing in the pipeline was developed on or changed for, it gave no
wrong answer, and almost no answers.

\section{Every wrong answer, and where it came from}
\label{sec:ledger}

Table~\ref{tab:ledger} lists the wrong answers recorded during the
pipeline's development and in both held-out tests. Each was traced to
one step of the translation, and each FraCaS one is now prevented by a
general rule rather than a per-problem case.

\begin{table}[t]
\centering\small
\setlength{\tabcolsep}{3pt}
\begin{tabular}{L{1.45cm}L{3.3cm}L{2.2cm}}
\toprule
problems & cause & now \\
\midrule
FraCaS 56, 77, 94, 95, 109 & committing to one reading (``many'' as ``some'', ``both'' as ``every'', \ldots) & every-reading rule \\
FraCaS 30, 32, 47 & one name for ``spend a lot of time'' and ``spend time'' & determiners in names; name check \\
FraCaS 304 & ``for two hours'' deleted from ``wrote a report'' & telicity \\
FraCaS 78 & a \no{} claim ignoring a partitive's presupposition & presupposition in \no{} claims \\
MED 2718 & one name for ``said he didn't drink black coffee'' and the phrase inside it & name check (post-hoc) \\
MED 536, 2504 & MED's swapped-pair construction & counted as wrong \\
SICK & none & \\
\bottomrule
\end{tabular}
\caption{The wrong answers recorded during development (FraCaS) and in
the held-out runs.}
\label{tab:ledger}
\end{table}

Two things stand out. First, the errors fall into a few classes, and
the two largest, committing to a reading and colliding names, each
have one general fix. Second, the held-out error of the pipeline's own
making belongs to a class already seen on FraCaS: MED~2718 is the same
collision as FraCaS 30, through a word class (verbs taking a clause)
that FraCaS never exposed. This is why the name check tests the
collision itself rather than listing the word classes a name must keep.

\section{Related work}
\label{sec:related}

\paragraph{Logic-based FraCaS systems.} ccg2lambda
\citep{mineshima-etal-2015-higher} composes higher-order formulas from
CCG parses and proves them in Coq with a 30-second timeout; its authors
note that their FraCaS figures are not an evaluation on unseen data.
\citet{haruta-etal-2020-logical} extend it to comparatives and
generalized quantifiers, and also evaluate on the MED subset above and
on a new set of comparative problems. FraCoq
\citep{bernardy-chatzikyriakidis-2017-type,bernardy-chatzikyriakidis-2019-wide,bernardy-chatzikyriakidis-2021-applied}
translates GF trees from the FraCaS treebank
\citep{ljunglof-siverbo-2011-bilingual} into Coq and covers all sections,
with partly human-built proofs; its authors stress that the value of
the work lies in how the inferences are made rather than in the bare
accuracy numbers. Our pipeline differs in parsing FraCaS's English,
proving without human help, and answering only when every reading
agrees. ccg2lambda's FraCaS results do not report wrong answers
separately from accuracy.

\paragraph{Precision on SICK and MED.} LangPro
\citep{abzianidze-2015-tableau,abzianidze-2017-langpro}, a natural-logic
tableau prover, and
\citet{martinez-gomez-etal-2017-demand}, ccg2lambda with lexical
knowledge injected on demand, are reported at 98.0 and 97.0 precision and
58.1 and 63.6 recall on SICK \citep[Table 3]{chen-etal-2021-neurallog}. NeuralLog \citep{chen-etal-2021-neurallog}
combines a monotonicity engine with neural phrase alignment and reaches
93.4 on MED. Our precision is not shown to be higher than LangPro's or
Mart\'inez-G\'omez et al.'s, and our recall is far lower. What is
different is the source of the guarantee: every answer is a checked
proof of a stated claim, so every error can be traced to a line of the
translation, as each of ours was.

\paragraph{Dataset construction.} MED's crowdsourced pairs are labelled
by edit direction and then swapped with reversed labels
\citep{yanaka-etal-2019-neural}. Two of our three MED errors are pairs
where that procedure went wrong; a sound system makes such pairs visible
as wrong answers, though it cannot tell them from its own errors without
inspection.

\section{Conclusion}

A logic-based NLI system with checked proofs can be evaluated, and
debugged, by its wrong answers alone: each one points to a step of the
translation. On FraCaS this led to two general rules, answer only when
every reading agrees, and refuse a name that would stand for two
phrases, which together with a few semantic safeguards remove every
wrong answer the pipeline gave. Held out, the pipeline made one
translation error on MED, of exactly the kind the name check now
refuses, and none on a pre-registered SICK test. The price is recall:
the system answers very little outside FraCaS. Raising recall without
giving up the guarantee means more grammar and lexical knowledge, each
of which is new translation, and each must again be tested on data it
was not developed on.

\section*{Limitations}

\begin{itemize}
\item \textbf{Development-set figures.} Every FraCaS construction,
reading, background rule and prover step was added after a FraCaS
failure; the name check was added after a MED failure. Only the SICK
run tests the pipeline as it is.
\item \textbf{Coverage.} The grammar is a controlled language fit to
FraCaS. It reaches 18.5\% of MED and 28.9\% of SICK, and recall on SICK
is 2.5\%. On Haruta et al.'s MED subset the pipeline is at the
all-\unk{} baseline.
\item \textbf{Small precision samples.} 52 held-out answers on SICK and
140 on MED cannot distinguish our precision from that of LangPro or
Mart\'inez-G\'omez et al.
\item \textbf{What the name check does not cover.} It compares phrases by
the words their names keep and drop, not by their structure, so two
phrases with the same words in a different structure would pass. A
construction with no word of its own is covered only when listed (the
progressive and the passive are). A false background premise is outside
it: the telicity ablation's two wrong answers are of that kind. The
deliberate omissions are simplifications: the claims are tenseless.
Two exemptions are not compared: a node the translator names as a whole
(``the items'' as \nm{every}), and one idiom named by part of its
phrase (``comes cheap'' as \nm{cheap}).
\item \textbf{Finite domains.} Claims quantify over an arbitrary
\emph{finite} domain, so an inference valid only in finite models would
be accepted.
\item \textbf{No lexical knowledge.} MED's and SICK's lexical problems
can only be \unk{}. Adding lexical knowledge would make some of the
collisions the name check refuses matter (MED~4897).
\item \textbf{Trusted base.} The guarantee assumes Lean's kernel and the
definitions of the generalized quantifiers are correct.
\item \textbf{Licences.} MED is CC BY-SA 4.0 and SICK CC BY-NC-SA 3.0;
the files derived from them carry those licences.
\end{itemize}

\bibliography{references}

\end{document}
